\section{Methodology}
\label{sec:methodology}

We introduce Instruction-Prefix-Conditioned Routing (IPCR), a minimal approach for zero-shot adapter selection. Our design philosophy is to test whether intrinsic alignment exists before adding architectural complexity.

\subsection{Overview}

Given an instruction $x$, IPCR routes to a bank of $k$ task-specific LoRA adapters $\{\theta_1, \ldots, \theta_k\}$ through three steps:

\begin{enumerate}
    \item \textbf{Embed:} Encode the instruction prefix using a frozen sentence encoder $f: \mathcal{X} \rightarrow \mathbb{R}^d$
    \item \textbf{Route:} Apply a linear probe $W \in \mathbb{R}^{k \times d}$ to predict adapter scores
    \item \textbf{Select:} Choose the highest-scoring adapter (hard routing) or compute weighted combination (soft routing)
\end{enumerate}

\textbf{Rationale:} If instruction semantics and adapter specializations share geometric structure, a linear mapping should suffice. This tests our hypothesis directly---any success demonstrates intrinsic alignment, and failure would motivate more complex architectures.

\subsection{Instruction Embedding}

We use MiniLM-L6-v2~\cite{wang2020minilm}, a 22M-parameter sentence encoder pretrained via contrastive learning, as our embedding function $f$. The encoder is completely frozen during both probe training and inference.

\textbf{Why MiniLM?} Three considerations:
\begin{enumerate}
    \item \textbf{Efficiency:} 22M parameters, 384-dimensional output, negligible latency overhead
    \item \textbf{Pretraining:} Contrastive objective should provide task-relevant semantic clusters
    \item \textbf{Simplicity:} Tests whether off-the-shelf embeddings contain routing signal
\end{enumerate}

We encode only the instruction prefix (first 128 tokens), excluding input/output content. This focuses routing on task semantics rather than instance-specific features.

\subsection{Linear Probe Training}

The routing probe is a multinomial logistic regression classifier:

\begin{equation}
p(y = j | x) = \frac{\exp(W_j \cdot f(x))}{\sum_{i=1}^k \exp(W_i \cdot f(x))}
\end{equation}

where $W_j$ is the weight vector for adapter $j$. Training uses L-BFGS optimization with balanced class weights to handle task-frequency imbalance in FLAN.

\textbf{Why linear?} A linear probe has minimal capacity---it can only exploit structure that already exists in the embedding space. If the probe succeeds, we have evidence that MiniLM embeddings encode task-relevant structure. A complex nonlinear router would conflate learned routing with intrinsic alignment.

\subsection{Adapter Bank}

Each adapter $\theta_j$ is a task-specialized LoRA~\cite{hu2021lora} with:
\begin{itemize}
    \item Rank $r = 16$
    \item Scaling factor $\alpha = 32$
    \item Applied to query/value projections in all attention layers
\end{itemize}

Adapters are trained independently on their respective task families, then frozen for routing experiments. This ensures the routing probe cannot modify adapter behavior---it can only select among fixed options.

\subsection{Routing Strategies}

\textbf{Hard routing:} Select adapter with highest probe score: $\hat{j} = \arg\max_j W_j \cdot f(x)$

\textbf{Soft routing:} Compute weighted combination of adapter outputs:
\begin{equation}
\Delta W = \sum_{j=1}^k \sigma(W_j \cdot f(x)) \cdot \theta_j
\end{equation}

where $\sigma$ is the softmax function. We report hard routing results in main experiments and analyze soft routing in ablations.

\subsection{Baseline Comparisons}

To isolate routing contribution, we compare against:
\begin{itemize}
    \item \textbf{Oracle:} Always select the adapter trained on the input's task family (upper bound)
    \item \textbf{Uniform:} Equal weight to all adapters (no routing information)
    \item \textbf{Random:} Uniformly random adapter selection
\end{itemize}

This setup tests whether IPCR extracts meaningful routing signal versus simple aggregation strategies.
